\documentclass[10pt,a4paper]{article}

\usepackage[utf8]{inputenc}
\usepackage[T1]{fontenc}
\usepackage{lmodern}
\usepackage[margin=1.5cm]{geometry}
\usepackage{microtype}
\usepackage{amsmath,amssymb}
\usepackage{enumitem}
\usepackage{titlesec}
\usepackage{xcolor}
\usepackage[hidelinks]{hyperref}

\definecolor{primary}{RGB}{20, 42, 92}
\definecolor{muted}{RGB}{90, 98, 112}

\pagestyle{empty}
\setlength{\parindent}{0pt}
\setlist[itemize]{leftmargin=1.4em, nosep, topsep=2pt, itemsep=1.5pt}

\titleformat{\section}{\large\bfseries\color{primary}\uppercase}{}{0em}{}[\titlerule]
\titlespacing*{\section}{0pt}{8pt}{4pt}

\begin{document}

% ==================== HEADER ====================
\begin{center}
  {\LARGE \textbf{\color{primary} Julian Leonardo Avila-Martinez}}\\[3pt]
  {\small Theoretical \& Computational Physicist $\cdot$ Scientific Software Engineer $\cdot$ AI-for-Science \& Mathematical Modeling}\\[3pt]
  {\footnotesize
    Bogot\'a D.C., Colombia $\cdot$
    \href{mailto:jlavilam@udistrital.edu.co}{\texttt{jlavilam@udistrital.edu.co}} $\cdot$
    \href{mailto:julianlavilam@gmail.com}{\texttt{julianlavilam@gmail.com}} $\cdot$
    \textbf{Available Immediately}
  }\\[2pt]
  {\footnotesize
    \textbf{ORCID:} \href{https://orcid.org/0009-0002-8630-1922}{\texttt{0009-0002-8630-1922}} $\cdot$
    \textbf{CvLAC:} \href{https://scienti.minciencias.gov.co/cvlac/visualizador/generarCurriculoCv.do?cod_rh=0002012647}{\texttt{0002012647}} $\cdot$
    \textbf{GitHub:} \href{https://github.com/Julian-L-Avila}{\texttt{github.com/Julian-L-Avila}} $\cdot$
    \textbf{Web:} \href{https://julian-l-avila.github.io/JulianLAvila/}{\texttt{julian-l-avila.github.io/JulianLAvila}}
  }
\end{center}

% ==================== EDUCATION ====================
\section{Education}

\textbf{Universidad Distrital Francisco Jos\'e de Caldas (UDFJC)} \hfill \textbf{Bogot\'a D.C., Colombia}\\
\textit{B.Sc. in Physics (Pregrado en F\'isica --- 5-Year Degree, 10 Semesters, 145 Credits / $\ge 240$ ECTS)} \hfill \textit{Nov 2021 -- Dec 2026}
\begin{itemize}
  \item \textbf{Academic Standing:} \textbf{Ranked 2nd in Graduating Cohort (Top 2\%)} $\cdot$ \textbf{Cumulative GPA: 4.54 / 5.00} (Cohort Max: 4.55 / 5.00; German Scale Equivalent: \textbf{1.7}; \textbf{ECTS Grade A}). All 145 credits and thesis defense completed (\textit{Acta No. 020 de 2026}); formal graduation ceremony December 2026 (\textbf{available immediately}).
  \item \textbf{Undergraduate Thesis (Monograph, Written in English):} \textit{``Foundations of the Dirac Equation: From Relativistic Principles to Clifford Algebras''}.
  \item \textbf{Thesis Advisor:} \textbf{Dr. Asdr\'ubal Moreno-Mosquera} (Ph.D. in Physics, Universidad de Salamanca, Spain --- IUFFyM, \texttt{amorenom@udistrital.edu.co}) $\cdot$ \textbf{Honors:} \textbf{Nominated for Meritorious Thesis (\textit{Menci\'on Meritoria})}.
\end{itemize}

% ==================== AWARDS AND HONORS ====================
\section{Awards and Honors}

\begin{itemize}
  \item \textbf{Academic Excellence Honor Distinction (\textit{Matr\'icula de Honor}), UDFJC (2021 -- 2025):} Awarded in \textbf{8 of 9 evaluated academic semesters} to the top 10 GPAs of the physics program with zero failed courses (semesters \texttt{2021-1}, \texttt{2022-3}, \texttt{2023-1}, \texttt{2023-3}, \texttt{2024-1}, \texttt{2024-3}, \texttt{2025-1}, and \texttt{2025-3}).
  \item \textbf{Meritorious Thesis Nomination (\textit{Menci\'on Meritoria}), UDFJC (2026):} Nominated for the undergraduate thesis monograph \textit{``Foundations of the Dirac Equation: From Relativistic Principles to Clifford Algebras''}.
  \item \textbf{Third Best Poster Award (Particle \& High-Energy Physics Category), CIFI-SEMA 2026 (May 2026):} Congreso Internacional de F\'isica e Ingenier\'ia, Universidad Sim\'on Bol\'ivar, Barranquilla (\textit{``Chiral Dirac Equation: Degrees of Freedom and Neutrino's Mass''}).
  \item \textbf{Best Poster Award (Research Seedbed Category), I Semana de las Ciencias, UDFJC (Nov 2023):} Facultad de Ciencias Matem\'aticas y Naturales (\textit{``Asymmetric Capacitor -- Lifter''}, CUEVA Research Group).
\end{itemize}

% ==================== RESEARCH & SCIENTIFIC COMPUTING EXPERIENCE ====================
\section{Research \& Scientific Computing Experience}

\textbf{Spinor Theory, Clifford Algebras, Chiral Fermions \& Relativistic Quantum Software} \hfill \textbf{2025 -- 2026}\\
\textit{Advisors: Dr. Asdr\'ubal Moreno-Mosquera (UDFJC) and Dr. Jos\'e Alfonso Leyva Rojas (PU Javeriana)}
\begin{itemize}
  \item Formulated the free and interacting Dirac field from first principles via the representation theory of real Clifford algebras $\mathcal{C}\ell_{1,3}(\mathbb{R})$, the double cover $\mathrm{Spin}^+(1,3) \to \mathrm{SO}^+(1,3)$, and the Yvon--Takabayashi polar decomposition $\psi = \rho^{1/2} e^{\beta I / 2} R$.
  \item Derived pseudo-Hermitian chiral mass generation $M = m_1 \mathbb{I}_4 + i m_2 \gamma^5$ ($\eta H \eta^{-1} = H^\dagger$), obtaining seesaw mass splittings $m_\pm = m_1 \pm \sqrt{m_1^2 - m_2^2}$ for neutrino phenomenology without auxiliary scalar fields.
  \item Co-developed \textbf{\texttt{dirac-solver}}, an open-source C++17 / \texttt{pybind11} / Python scientific library (MIT License; software registration pending) implementing staggered-leapfrog FDTD and Split-Operator FFT integrators for the (1D/2D/3D) Dirac equation with Builder/Strategy design patterns, Doxygen documentation, HDF5 I/O, and \texttt{pytest} validation suites.
\end{itemize}

\vspace{3pt}
\textbf{CUEVA Research Group, UDFJC --- Researcher in Relativistic Magnetohydrodynamics (RRMHD)} \hfill \textbf{2024 -- Present}\\
\textit{Advisor: Dr. Sergio Miranda Aranguren (Ph.D., Universitat de Val\`encia)}
\begin{itemize}
  \item Derived $3+1$ covariant tensor representations of Resistive Relativistic Magnetohydrodynamics (RRMHD), coupling stress-energy conservation ($\nabla_\mu T^{\mu\nu} = 0$) with Maxwell's equations and the relativistic Ohm's law.
  \item Modeled nonlinear hyperbolic conservation laws governing relativistic astrophysical slab jets and resistive current sheets; implemented numerical simulation pipelines in \textbf{Fortran 95 and Python} to benchmark jet collimation and resistive instabilities.
\end{itemize}

\vspace{3pt}
\textbf{Computational Physics, Scientific Machine Learning (SciML) \& AI Benchmark Engineering} \hfill \textbf{2024 -- 2026}\\
\textit{Department of Physics, UDFJC}
\begin{itemize}
  \item \textbf{Numerical PDEs, ODEs \& HPC Benchmarking:} Implemented and validated $\mathcal{O}(h^6)$ Numerov quantum solvers, graphene tight-binding dispersion, Velocity-Verlet Fortran 90 Molecular Dynamics, C++ 2D pipe fluid flow, Julia double-pendulum dynamics, and a 5-language 2D Laplace PDE benchmark (Fortran 90: \texttt{0.097 s}, C: \texttt{0.111 s}, MATLAB: \texttt{0.130 s}, C++: \texttt{0.172 s}, Python: \texttt{7.163 s}; $\sim 73.8\times$ compiled speedup).
  \item \textbf{SciML, PINNs \& High-Energy Data Mining:} Built Physics-Informed Neural Networks (PINNs) and 2D Galerkin Finite Element Method (FEM) solvers; benchmarked 11 ML architectures on 36,358 LHC collision events, achieving \textbf{90.3\% accuracy (0.958 ROC-AUC)} with a 2-layer MLP (\texttt{32-16}, ReLU).
  \item \textbf{AI-for-Science Benchmark Problems \& Automated Graders:} Created multi-stage quantum mechanics benchmark tasks, rubrics, and step-by-step \LaTeX{} reference derivations (4D Hilbert space CSCO $\{H,A,B\}$) to audit LLM symbolic reasoning; authored automated \texttt{pytest}/\texttt{unittest} test suites and C++ Integral Absolute Error (IAE) model validators against analytical and experimental data.
\end{itemize}

% ==================== PUBLICATIONS, SOFTWARE & PRESENTATIONS ====================
\section{Publications, Scientific Software \& Conference Presentations}

\begin{itemize}
  \item \textbf{[Monograph]} \textbf{Avila-Martinez, J. L.} (2026). \textit{Foundations of the Dirac Equation: From Relativistic Principles to Clifford Algebras}. B.Sc. Thesis Monograph, UDFJC (\textbf{Nominated for \textit{Menci\'on Meritoria}}).
  \item \textbf{[Working Paper]} \textbf{Avila-Martinez, J. L.}, Herrera-Martinez, L. Y., \& Rodriguez-Garcia, S. (2026). \textit{Hermiticity and Mass Generation in the Chiral Dirac Equation} (Working paper in preparation).
  \item \textbf{[Software]} \textbf{Avila-Martinez, J. L.}, Rodriguez-Garcia, S., \& Huertas-Archila, C. A. (2026). \textit{\texttt{dirac-solver}: Scientific Python Library with C++17 Backend for the Numerical Integration of the Dirac Equation} (Software registration pending).
  \item \textbf{[Talk]} \textit{The Vector Algebra War: Fragmentation in the Mathematical Formalisms of Physics}, Simposio de Historia de la Ciencia, III Semana de la Ciencia, UDFJC, Bogot\'a (Oct 2026).
  \item \textbf{[Poster --- Award Winner]} \textit{Chiral Dirac Equation: Degrees of Freedom and Neutrino's Mass}, CIFI-SEMA 2026, Barranquilla (May 2026, Poster \texttt{FP6-P3} --- \textbf{3rd Best Poster Award in Particle Physics}).
  \item \textbf{[Talk]} \textit{Survey on the Classification of Topological Phases of Matter}, International Symposium on Quantum Science and Engineering (ISQSE 2025), Tunja (Nov 2025).
  \item \textbf{[Talk]} \textit{The Eikonal as a Bridge to the Schr\"odinger Equation: Classical Limit and JWKB Approximation}, XXX Congreso Nacional de F\'isica, Manizales (Oct 2025).
  \item \textbf{[Poster]} \textit{Estudio Num\'erico de Jets Relativistas en el Marco de la RRMHD}, XXX Congreso Nacional de F\'isica, Manizales (Oct 2025).
  \item \textbf{[Poster]} \textit{Slab Jet: Numerical Study Using the RRMHD Framework}, Colombian Congress of Astronomy and Astrophysics (CoCoA 2024), Bucaramanga (Nov 2024).
  \item \textbf{[Poster --- Award Winner]} \textit{Asymmetric Capacitor -- Lifter}, I Semana de las Ciencias, UDFJC, Bogot\'a (Nov 2023 --- \textbf{Best Poster Award in Research Seedbed Category}).
\end{itemize}

% ==================== TEACHING, LEADERSHIP & ACADEMIC SERVICE ====================
\section{Teaching, Governance \& Scientific Event Organization}

\begin{itemize}
  \item \textbf{Founder \& Seminar Leader --- Clifford Seminar Group, UDFJC (2026 -- Present):} Founded and guides the weekly research seminar on Clifford algebras and spinor geometry using the German \textit{Thesenpapier} methodology.
  \item \textbf{Academic Monitor (Teaching Assistant), Dept. of Physics, UDFJC (Terms 2025-1, 2026-1):} Led 14 hrs/week of recitation, problem-set/rubric design, and reference solution authoring in \textit{Mathematical Methods for Physicists} (Lie algebras, quaternions, Sturm--Liouville, Lorentz group, PDEs), \textit{Vector Calculus}, \textit{Linear Algebra}, and \textit{Physics I \& II}.
  \item \textbf{Founding Member --- Colombian Physics Society (\textit{Acolf\'isica}) Student Chapter at UDFJC (2025 -- Present).}
  \item \textbf{Elected Student Representative --- Curricular Council of Physics (\textit{Consejo Curricular de F\'isica}), UDFJC (2024 -- 2025).}
  \item \textbf{Scientific Event Co-Organizer (UDFJC -- PU Javeriana):} Co-organized \textit{Simposio Centenario de la Ecuaci\'on de Schr\"odinger} (2026, \href{https://udpuj.github.io/simposio-centenario-schr-dinger}{\texttt{udpuj.github.io/simposio-centenario-schr-dinger}}) and \textit{Simposio La Mec\'anica Cu\'antica en su Centenario} (Nov 2025, \href{https://udpuj.github.io/simposio-centenario-cuantica/}{\texttt{udpuj.github.io/simposio-centenario-cuantica/}}); committee member for \textit{Horizontes Cu\'anticos} (Sep 2025) and \textit{II Semana de la Ciencia} (Oct 2024).
\end{itemize}

% ==================== TECHNICAL SKILLS & LANGUAGES ====================
\section{Technical Stack, Languages \& Academic References}

\begin{itemize}
  \item \textbf{Scientific Python \& SciML Stack:} \texttt{NumPy}, \texttt{SciPy}, \texttt{SymPy}, \texttt{PyTorch} (PINNs/Deep Learning), \texttt{scikit-learn}, \texttt{Pandas}, \texttt{h5py} (HDF5), \texttt{pybind11}, \texttt{Matplotlib}, \texttt{pytest}, \texttt{unittest}, Jupyter Notebooks.
  \item \textbf{Compiled Languages \& Engineering Tools:} \textbf{C / C++17} (CMake, OOP, FDTD/ODEs), \textbf{Fortran 90/95} (Molecular Dynamics, RRMHD, PDEs), \textbf{Julia}, \textbf{MATLAB}, GNU/Linux (Debian), Git, Bash, Doxygen, \LaTeX{} (Bib\LaTeX, Beamer, TikZ, \texttt{circuitikz}), Gnuplot, KiCad.
  \item \textbf{Languages:} \textbf{Spanish} (Native) $\cdot$ \textbf{English} (\textbf{C1 CEFR --- IELTS Academic 8.0}: Listening 9.0, Reading 8.5, Speaking 7.5, Writing 6.0; thesis written in English) $\cdot$ \textbf{German} (\textbf{B1 CEFR}).
  \item \textbf{Academic References:}
    \textbf{Dr. Asdr\'ubal Moreno-Mosquera} (Tenured Prof., UDFJC; Ph.D. U. de Salamanca, \texttt{amorenom@udistrital.edu.co}) $\cdot$
    \textbf{Dr. Jos\'e Alfonso Leyva Rojas} (Tenured Prof., PU Javeriana; Ph.D. JGU Mainz, \texttt{leyvaa@javeriana.edu.co}) $\cdot$
    \textbf{Dr. Sergio Miranda Aranguren} (Assoc. Prof., UDFJC; Ph.D. U. de Val\`encia, \texttt{smirandaa@udistrital.edu.co}).
\end{itemize}

\end{document}
